
\subsubsection{Research Questions}

\textbf{RQ1 (CodeBLEU Reliability):} Does structural similarity demonstrate measurement reliability across all three criteria?

\textbf{RQ2 (Runtime Efficiency Threshold):} Does runtime efficiency measurement achieve the CV $\leq 5$\% threshold?

\textbf{RQ3 (Gate Logic Functionality):} Does the multi-criteria validation gate correctly discriminate between reliable and unreliable proxies?

\subsubsection{Proof-of-Concept Scope}

This Phase 4 implementation is a proof-of-concept validating methodology before infrastructure investment. Full implementation would require: GPU infrastructure (CodeLlama-7B-Instruct, 16GB+ VRAM), HumanEval dataset (164 problems), hardware performance tools (Linux \texttt{perf}), and cross-platform setup (AWS + local GPU).

\textbf{PoC Strategy:} We validate the statistical framework, gate logic, and visualization pipeline using synthetic measurements (500 solutions $\times$ 5 repetitions $\times$ 3 metrics = 7,500 data points).

\textbf{PoC Success Criteria:} Code executes without error, statistical computations produce interpretable results, gate logic discriminates between passing and failing proxies, and at least one proxy validates.

\subsubsection{Measurement Reliability Criteria}

\textbf{Coefficient of Variation (CV $\leq 5$\%):} Measures intra-implementation variability. Threshold ensures 95\% of variance reflects code quality differences, not measurement noise.

\textbf{Cohen's d ($\geq 0.8)$:} Measures inter-group separation between complexity classes. Threshold ensures proxies distinguish efficient from inefficient implementations.

\textbf{Spearman $\rho$ ($\geq 0.8)$:} Measures cross-platform rank preservation. Threshold ensures measurements are platform-invariant.
